\documentclass[12pt,a4paper]{article}
\usepackage[margin=25mm, headheight=15pt, headsep=10pt]{geometry}
\usepackage{fontspec}
\usepackage[table]{xcolor} % Note: table option is required for rowcolors
\usepackage{titlesec}
\usepackage{tocloft}
\usepackage{fancyhdr}
\usepackage{booktabs}
\usepackage{tcolorbox}
\tcbuselibrary{skins, breakable}
\usepackage[colorlinks=true, linkcolor=emerald, urlcolor=emerald, bookmarks=true]{hyperref}

\definecolor{emerald}{RGB}{0,102,51}
\definecolor{darkred}{RGB}{139,0,0}
\definecolor{lightgray}{RGB}{245,245,245}

\usepackage[nil]{babel}
\babelprovide[import, main]{bengali}
\babelprovide[import]{arabic}
\babelfont{rm}[Renderer=HarfBuzz,Scale=1.0]{Noto Serif Bengali}
\babelfont{sf}[Renderer=HarfBuzz]{Noto Sans Bengali}
\babelfont{tt}[Renderer=HarfBuzz]{Noto Sans Bengali}
\babelfont[arabic]{rm}[Renderer=HarfBuzz,Scale=1.1]{Noto Naskh Arabic}

% Section styling
\titleformat{\section}{\Large\bfseries\color{emerald}}{\thesection.}{0.5em}{}
\titleformat{\subsection}{\large\bfseries\color{emerald!85!black}}{\thesubsection.}{0.5em}{}
\titleformat{\subsubsection}{\normalsize\bfseries\color{emerald!70!black}}{\thesubsubsection.}{0.5em}{}
\titlespacing*{\section}{0pt}{18pt}{8pt}

% Fancy Header/Footer setup
\pagestyle{fancy}
\fancyhf{} % clear all header and footer fields
\fancyhead[L]{\color{emerald}\small\textbf{\nouppercase{\leftmark}}}
\fancyfoot[L]{\scriptsize\color{black!50}{\lat{AI}}-সংকলিত, আলেম-যাচাইকৃত নয়}
\fancyfoot[C]{\color{emerald!70!black}\thepage}
\renewcommand{\headrulewidth}{0.4pt}
\renewcommand{\headrule}{\hbox to\headwidth{\color{emerald}\leaders\hrule height \headrulewidth\hfill}}
\renewcommand{\footrulewidth}{0pt}

% Custom boxes
% 1. Ayat/Hadith Box
\newtcolorbox{ayatbox}[1][]{
    enhanced, breakable, 
    colback=emerald!5, 
    colframe=emerald, 
    boxrule=0.5pt, 
    arc=3pt, 
    left=10pt, right=10pt, top=10pt, bottom=10pt,
    #1
}

% 2. Quote/Translation Box
\newtcolorbox{quotebox}[1][]{
    enhanced, breakable,
    frame hidden,
    colback=lightgray,
    borderline west={3pt}{0pt}{emerald},
    left=15pt, right=10pt, top=10pt, bottom=10pt,
    sharp corners,
    #1
}

% 3. AI-generated notice box
\newtcolorbox{warnbox}[1][]{
    enhanced, breakable,
    colback=red!4,
    colframe=darkred,
    boxrule=0.8pt,
    arc=3pt,
    left=10pt, right=10pt, top=10pt, bottom=10pt,
    #1
}

% Commands
\newcommand{\arab}[1]{\foreignlanguage{arabic}{#1}}
\newcommand{\kib}[1]{\textcolor{emerald}{\textbf{#1}}}

% Latin (English) fallback: Noto Serif Bengali has NO Latin glyphs
\newfontfamily\latfont[Renderer=HarfBuzz]{Noto Serif}
\newcommand{\lat}[1]{{\latfont #1}}

\setlength{\parskip}{5pt}
\setlength{\parindent}{0pt}

% Fix for missing bullet in Noto Serif Bengali
\renewcommand{\labelitemi}{$\bullet$}



\begin{document}
\begin{center}{\Large\bfseries\color{emerald} \lat{11}-তাশাহহুদ-শেষ-দরুদ}\end{center}
\vspace{1em}
\section*{১১ — তাশাহহুদ (শেষ) ও দরুদ ইবরাহিম}

\begin{warnbox}
 \textbf{গুরুত্বপূর্ণ নোটিশ:} এই কন্টেন্টটি \lat{AI} দিয়ে প্রস্তুত। কেবল সহীহ/হাসান হাদিস রাখা হয়েছে। আলেম-যাচাইকৃত নয়।
\end{warnbox}

\medskip\hrule\medskip

\subsection*{১. সূত্র কার্ড}

\begin{table}[h]\centering\small
\begin{tabular}{ll}
বিষয় & বিবরণ \\
\hline
\textbf{বিষয়} & শেষ বৈঠকে পূর্ণ তাশাহহুদ + দরুদ ইবরাহিম \\
\textbf{সূত্র} & সহীহ বুখারি ৩৩৭০, সহীহ মুসলিম ৪০৬ \\
\textbf{মর্যাদা} & \textbf{ফরজ} (তাশাহহুদ) ও \textbf{সুন্নাতে মুআক্কাদাহ} (দরুদ) \\
\hline
\end{tabular}
\end{table}

\medskip\hrule\medskip

\subsection*{২. দরুদ ইবরাহিম — মূল আরবি}

\begin{center}\arab{اللَّهُمَّ صَلِّ عَلَىٰ مُحَمَّدٍ وَعَلَىٰ آلِ مُحَمَّدٍ كَمَا صَلَّيْتَ عَلَىٰ إِبْرَاهِيمَ وَعَلَىٰ آلِ إِبْرَاهِيمَ إِنَّكَ حَمِيدٌ مَجِيدٌ، اللَّهُمَّ بَارِكْ عَلَىٰ مُحَمَّدٍ وَعَلَىٰ آلِ مُحَمَّدٍ كَمَا بَارَكْتَ عَلَىٰ إِبْرَاهِيمَ وَعَلَىٰ آلِ إِبْرَاهِيمَ إِنَّكَ حَمِيدٌ مَجِيدٌ}\end{center}

\medskip\hrule\medskip

\subsection*{৩. বাংলা উচ্চারণ}

\textbf{আল্লাহুম্মা সাল্লি আলা মুহাম্মাদিন ওয়া আলা আলী মুহাম্মাদিন, কামা সাল্লাইতা আলা ইবরাহীমা ওয়া আলা আলী ইবরাহীম, ইন্নাকা হামিদুন মাজীদ। আল্লাহুম্মা বারিক আলা মুহাম্মাদিন ওয়া আলা আলী মুহাম্মাদিন, কামা বারাকতা আলা ইবরাহীমা ওয়া আলা আলী ইবরাহীম, ইন্নাকা হামিদুন মাজীদ।}

\medskip\hrule\medskip

\subsection*{৪. শব্দ-শব্দ অর্থ}

\subsubsection*{\textbf{দরুদ ইবরাহিম:}}

\begin{table}[h]\centering\small
\begin{tabular}{lll}
আরবি & বাংলা & \lat{English} \\
\hline
\textbf{\arab{اللَّهُمَّ}} & হে আল্লাহ & \textit{\lat{O} \lat{Allah}} \\
\textbf{\arab{صَلِّ}} & রহমাত দাও/দরুদ পাঠাও & \textit{\lat{send} \lat{blessings}} \\
\textbf{\arab{عَلَىٰ} \arab{مُحَمَّدٍ}} & মুহাম্মদের ওপর & \textit{\lat{upon} \lat{Muhammad}} \\
\textbf{\arab{وَعَلَىٰ} \arab{آلِ} \arab{مُحَمَّدٍ}} & ও মুহাম্মদের বংশের ওপর & \textit{\lat{and} \lat{upon} \lat{the} \lat{family} \lat{of} \lat{Muhammad}} \\
\textbf{\arab{كَمَا} \arab{صَلَّيْتَ}} & যেমন তুমি পাঠিয়েছ & \textit{\lat{as} \lat{You} \lat{sent}} \\
\textbf{\arab{عَلَىٰ} \arab{إِبْرَاهِيمَ}} & ইবরাহিমের ওপর & \textit{\lat{upon} \lat{Ibrahim}} \\
\textbf{\arab{إِنَّكَ}} & নিশ্চয়ই তুমি & \textit{\lat{Indeed} \lat{You} \lat{are}} \\
\textbf{\arab{حَمِيدٌ}} & প্রশংসিত & \textit{\lat{Praiseworthy}} \\
\textbf{\arab{مَجِيدٌ}} & মহিমান্বিত & \textit{\lat{Glorious}} \\
\textbf{\arab{بَارِ}�\arab{ْ}} & বরকত দাও & \textit{\lat{bless}} \\
\textbf{\arab{كَمَا} \arab{بَارَكْتَ}} & যেমন তুমি বরকত দিয়েছ & \textit{\lat{as} \lat{You} \lat{blessed}} \\
\hline
\end{tabular}
\end{table}

\medskip\hrule\medskip

\subsection*{৫. পূর্ণ বাংলা অনুবাদ}

\textbf{\lat{“}হে আল্লাহ! মুহাম্মদের ও মুহাম্মদের বংশের ওপর রহমাত পাঠাও — যেমন তুমি ইবরাহিমের ও ইবরাহিমের বংশের ওপর রহমাত পাঠিয়েছ। নিশ্চয়ই তুমি প্রশংসিত, মহিমান্বিত।}

\textbf{হে আল্লাহ! মুহাম্মদের ও মুহাম্মদের বংশের ওপর বরকত দাও — যেমন তুমি ইবরাহিমের ও ইবরাহিমের বংশের ওপর বরকত দিয়েছ। নিশ্চয়ই তুমি প্রশংসিত, মহিমান্বিত।\lat{”}}

\medskip\hrule\medskip

\subsection*{৬. সংক্ষিপ্ত ব্যাখ্যা}

\subsubsection*{\textbf{দরুদ ইবরাহিম কী?}}
নবীজি \arab{ﷺ}-র ও তাঁর বংশের ওপর আল্লাহর রহমাত ও বরকত প্রার্থনা। নবীজি \arab{ﷺ} থেকে ইবরাহিম \arab{ﷺ}-এর সাথে সম্পর্ক স্থাপন করা হয়েছে — কারণ দুজনই আল্লাহর হাবীব।

\subsubsection*{\textbf{"সাল্লি" ও "বারিক" — দুটো ভিন্ন প্রার্থনা}}
\textbf{"\lat{Salli} (সাল্লি)"} — রহমাত/দয়া প্রার্থনা। \textbf{"\lat{Barik} (বারিক)"} — বরকত/দীর্ঘস্থায়ী কল্যাণ প্রার্থনা। দুটো মিলিয়ে \textbf{\lat{complete} \lat{blessing} (সম্পূর্ণ বরকত)}।

\subsubsection*{\textbf{"আলী" — বংশ ও পরিবার}}
"\lat{Al}" মানে পরিবার/বংশ। মুহাম্মদ \arab{ﷺ}-এর পরিবার হলো বনূ হাশিম — যাদের মধ্যে ফাতিমা, আলী, হাসান, হুসাইন অন্যতম।

\medskip\hrule\medskip

\subsection*{৭. খুশুর জন্য মূল পয়েন্টস}

\begin{table}[h]\centering\small
\begin{tabular}{ll}
পয়েন্ট & কীভাবে প্রয়োগ করবেন \\
\hline
\textbf{নবীজির স্মরণ} & দরুদ পড়ার সময় নবীজি \arab{ﷺ}-কে মনে করুন — তিনি আমাদের জন্য কী করেছেন \\
\textbf{ধীরতা} & দরুদ ধীরে, স্পষ্ট করে পড়ুন — তাড়াহুড়ো নয় \\
\textbf{পুনরাবৃত্তি} & কমপক্ষে ১ বার; উত্তম ৩ বার \\
\textbf{দোয়া} & দরুদের পর যা চান তা চাইতে পারেন — দরুদের পর দোয়া কবুল হওয়ার সম্ভাবনা বেশি \\
\hline
\end{tabular}
\end{table}

\medskip\hrule\medskip

\subsection*{৮. সংশ্লিষ্ট হাদিস}

\begin{table}[h]\centering\small
\begin{tabular}{ll}
হাদিস & গ্রন্থ \\
\hline
\textbf{\lat{“}যে ব্যক্তি আমার ওপর একবার দরুদ পাঠায়, আল্লাহ তার ওপর ১০ বার রহমাত পাঠান।\lat{”}} & সহীহ মুসলিম ৩৮৪ \\
\textbf{\lat{“}দরুদ পড়ো এবং দোয়া করো — কারণ দরুদের পর দোয়া ফেরে দেওয়া হয় না।\lat{”}} & মুসনাদে আহমাদ ১১০১০ (সহীহ) \\
\hline
\end{tabular}
\end{table}

\medskip\hrule\medskip

\textit{শেষ আপডেট: আগস্ট ২০২৬}

\end{document}
